\documentclass[11pt]{article}

\usepackage[margin=1in]{geometry}
\usepackage{graphicx}
\usepackage{booktabs}
\usepackage{hyperref}
\usepackage{siunitx}
\usepackage{float}
\usepackage{xurl}

\hypersetup{colorlinks=true, linkcolor=blue, citecolor=blue, urlcolor=blue}

\title{Physically Motivated Time-Series Photometry of AM CVn}
\author{TODO: Team Members}
\date{TODO: Submission Date}

\begin{document}
\maketitle

\begin{abstract}
We observed AM Canum Venaticorum (AM CVn) in the V band on 2026 May 20 to test
whether short-baseline CCD photometry can recover and refine its known
ultracompact variability.  We reduced 140 images with bias, dark, and
star-masked flat-field calibration, then performed AAVSO-sequence differential
photometry.  After rejecting two artifact frames, 138 exposures spanning
111.20 min remained.  A generic single-sine search detects a 525.06 s signal,
close to the expected first harmonic of AM CVn's positive superhump.  A
physically motivated double-wave model gives the best independent period,
\(1055.21 \pm 8.58\) s after empirical error scaling.  A fixed 1051.2 s
superhump-plus-harmonic template is nearly as good, but does not beat the free
model in BIC or cross-validation.  The data therefore support detection of AM
CVn's known superhump/harmonic family, but they do not tighten the literature
period constraints.
\end{abstract}

\section{Introduction}

AM CVn is the prototype of a class of hydrogen-deficient ultracompact
interacting binaries.  These systems contain an accreting white dwarf and a
helium-rich donor, and their periods are short enough to connect accretion-disc
physics, compact-binary evolution, and low-frequency gravitational-wave
astrophysics \cite{solheim2010,roelofs2006}.

The period structure of AM CVn is subtle.  Early photometry showed a double
wave near 17.5 min and a strong 8.76 min harmonic \cite{patterson1979}.  Later
spectroscopy established \(P_{\rm orb}=1028.7322\) s as the orbital period,
while the dominant \(1051.2\) s photometric signal is interpreted as a positive
superhump \cite{nelemans2001,roelofs2006}.  Recent work uses this period
structure to revisit AM CVn system parameters and period change
\cite{smak2023}.

Our goal is not only to detect variability, but to ask a sharper question: can
our local 111 min light curve constrain the literature periods more tightly?
To avoid overclaiming, we separate independent constraints from conditional
constraints.  Independent constraints use periods fitted freely from our data.
Conditional constraints assume a literature-motivated physical identification,
such as the known superhump or orbital signal.

\section{Observations and Reduction}

The observations were taken on 2026 May 20 in the V band.  The AM CVn sequence
contains 140 frames with exposure times of 35, 40, and 45 s.  The reduction used
late-night bias frames, exposure-matched darks, and a star-masked
\texttt{flat\_4s\_masked\_clipmean} flat.  The final two frames showed severe
artifacts and low counts, so they were excluded.

The science frames did not have a reliable WCS solution.  We therefore used the
AAVSO comparison sequence X42421BZ to map the target and comparison stars onto
a good reference frame.  Aperture photometry used a 6 px aperture and a
12--20 px sky annulus.  Fluxes were divided by exposure time before comparison,
so frames with different exposure times could be combined.  Each comparison
star was normalized by its own median flux before constructing the ensemble,
preventing the brightest comparison stars from dominating the differential
light curve.

The final light curve contains 138 good points.  It spans 111.20 min, or 6.35
cycles of a 1051.2 s superhump.  The robust target scatter is 1.51 percent,
while the median formal photometric error is 0.90 percent.  This excess scatter
motivates the empirical error scaling used below.

\section{Analysis Methods}

We first fit a model ladder.  The baseline model is constant flux.  The
generic models are a free single sine and a free double-wave model containing a
fundamental plus its first harmonic.  We also ran multi-harmonic
Lomb-Scargle periodograms: \texttt{nterms=2} is the periodogram analogue of the
free double-wave model.  The physically motivated models use literature
periods: the 1028.7322 s orbital period, the 1051.2 s positive superhump, a
1011.4 s negative superhump, and a superhump-plus-harmonic template.  Finally,
we fit a beat-constrained precessing-disc proxy with
\(f_{\rm sh}=f_{\rm orb}-f_{\rm prec}\).  The 13.38 h disc precession period
was not directly fit from the light curve because our run covers only 0.14
cycle of it.

For each model we report chi-square, BIC, RMS residual scatter, and
time-blocked cross-validation error.  For free periods, we compute profile
\(\chi^2\) intervals and block-bootstrap uncertainties.  Formal photon-noise
errors underestimate the residual scatter, so the final period uncertainties
use the best double-wave model's \(\sqrt{\chi^2_\nu}=1.59\) scale factor.  We
also test stability by fitting the first and second halves separately, clipping
outliers, dropping comparison stars one at a time, and checking residual
periodograms.  As final stress tests, we fit extra nuisance terms for time,
tracking shifts, exposure state, and ensemble common mode; repeat the
double-wave fit with robust loss functions; and rebuild the differential light
curve from comparison-star subsets.

\section{Results}

The single-sine search recovers \(525.06 \pm 4.27\) s after empirical error
scaling.  Ordinary Lomb-Scargle detects the same harmonic at 525.08 s.  When
Lomb-Scargle is generalized to a two-term Fourier model, it recovers a
fundamental near 1053.27 s.  Thus Lomb-Scargle is useful as a detection tool,
but it must be made multi-harmonic or followed by physical model fitting for AM
CVn.

The best independent model is the free double-wave fit:
\[
P_{\rm free}=1055.21 \pm 8.58~{\rm s}.
\]
This is consistent with the known 1051.2 s superhump family.  It is not precise
enough to improve the published photometric or spectroscopic periods.

\begin{table}[H]
\centering
\caption{AM CVn model ladder.  \(\Delta{\rm BIC}\) is measured relative to the
best free double-wave model; positive values are worse.}
\label{tab:model-ladder}
\resizebox{\textwidth}{!}{%
\begin{tabular}{llrrrr}
\toprule
Model & Physical interpretation & Period (s) & RMS (\%) & \(\Delta{\rm BIC}\) & CV RMS (\%) \\
\midrule
free double wave & fundamental + harmonic & 1055.21 & 1.43 & 0.00 & 1.52 \\
beat-constrained disk & \(f_{\rm sh}=f_{\rm orb}-f_{\rm prec}\) & 1055.22 & 1.43 & 0.00 & 1.52 \\
1051.2 s + harmonic & positive-superhump template & 1051.2 & 1.43 & 0.55 & 1.53 \\
free single sine & harmonic-only detection & 525.06 & 1.52 & 37.44 & 1.57 \\
orbital + positive + negative SH & three known clocks & fixed & 1.50 & 45.55 & 1.65 \\
orbital + positive SH & two known clocks & fixed & 1.53 & 52.99 & 1.71 \\
positive SH only & 1051.2 s sine & 1051.2 & 1.56 & 57.67 & 1.62 \\
orbital only & 1028.7322 s sine & 1028.73 & 1.58 & 66.18 & 1.65 \\
\bottomrule
\end{tabular}
}
\end{table}

The most important result in Table~\ref{tab:model-ladder} is that the
beat-constrained precessing-disc proxy and the free double-wave model are
effectively tied.  This happens because a freely fitted precession period maps
the known orbital period onto the best photometric double-wave period; the
derived value is \(P_{\rm prec}\simeq 11.39\) h.  This captures the expected
superhump physics better than a generic Fourier fit, but it is not an
independent precession detection because the baseline covers far less than one
precession cycle.  The fixed 1051.2 s superhump-plus-harmonic template is also
nearly tied, supporting the literature interpretation.

\begin{figure}[H]
\centering
\includegraphics[width=\textwidth]{analysis/project/plots/AM_CVn_physical_model_ladder.png}
\caption{Physically motivated model-suite summary.  The upper-left panel shows
the AM CVn differential light curve and the best free double-wave model.  The
upper-right panel folds the same model.  The lower-left panel shows that the
residual periodogram has no significant remaining peak.  The lower-right panel
compares BIC values against the free double-wave fit.}
\label{fig:model-suite}
\end{figure}

The validation tests are mixed.  Dropping one comparison star at a time changes
the free double-wave period by only 5.2 s, and clipping \(|z|>4\) outliers
changes it by only 0.02 s.  The residual periodogram has false-alarm
probability 0.99, so no strong unmodeled short-period signal remains.  However,
the first-half and second-half fitted periods differ by 22.8 s.  That split is
not surprising with only about six superhump cycles, but it blocks any strong
claim that the period is tightly constrained by our data alone.

We also tried three ways to squeeze more information out of the same light
curve.  First, adding nuisance regressors for time, image shifts, exposure
time, number of comparison stars, and common-mode behavior can reduce the raw
RMS from 1.43\% to 1.41\%, but BIC and blocked cross-validation get worse
(\(\Delta{\rm BIC}=22.66\), \(\Delta{\rm CV}=0.97\%\) for the most flexible
case).  This indicates overfitting rather than a better physical constraint.
Second, robust likelihoods do not sharpen the answer: soft-\(L_1\), Huber, and
Cauchy losses shift the best period to 1057--1059 s and increase the ordinary
RMS.  Third, selected comparison-star subsets can shave the target RMS as low
as 1.40\%, but the fitted period moves by a few seconds and the control-star
scatter does not justify replacing the full ensemble.  These tests are useful
because they show what does not work: the limiting uncertainty is the short
baseline and cycle-to-cycle structure, not just Gaussian photometric noise.

The broadband photometry cannot test AM CVn's S-wave emission directly.  The
S-wave is a spectroscopic velocity feature from the bright spot or stream-disc
impact region.  Testing it requires time-resolved spectroscopy or Doppler
tomography, not V-band imaging.  For the same reason, our photometry can support
the known white-dwarf-binary interpretation but cannot independently prove or
disprove the object class.

\begin{figure}[H]
\centering
\includegraphics[width=0.78\textwidth]{analysis/project/plots/AM_CVn_physical_literature_constraints.png}
\caption{Comparison between our period constraints and literature periods.  Our
independent scaled uncertainty is about 8.6 s, far larger than the published
precision of the orbital and superhump periods.  The result is therefore
consistent with the literature, not more constraining than it.}
\label{fig:literature-constraints}
\end{figure}

\section{Conclusions}

Our data clearly detect short-timescale AM CVn variability.  A single-sine
model finds a 525 s signal, while the physically preferred interpretation is a
double-wave signal in the known 1051 s superhump family.  The best independent
period is \(1055.21 \pm 8.58\) s, and the literature-fixed 1051.2 s
superhump-plus-harmonic template is nearly as good.

The result does not tighten the literature constraints.  The physically
motivated model does not beat the free double-wave model, the first/second-half
stability test fails, and our independent uncertainty is hundreds of times
larger than the best published period precision.  Detrending, robust
likelihoods, and comparison-star subset searches improve neither predictive
performance nor period stability.  The strongest defensible claim is therefore:
our photometry independently recovers variability consistent with AM CVn's
known superhump/harmonic structure, but a longer multi-night baseline would be
needed to improve published period constraints.

\section*{Acknowledgments}

We thank the PHYSICS 100 teaching staff for telescope access and reduction
guidance.  This work used AAVSO comparison-star resources and NASA ADS/arXiv
literature records.

\begin{thebibliography}{9}

\bibitem{nelemans2001}
Nelemans, G., Steeghs, D., \& Groot, P. J. 2001.
\textit{Spectroscopic evidence for the binary nature of AM CVn}.
MNRAS, 326, 621.
\url{https://academic.oup.com/mnras/article/326/2/621/1006901}.

\bibitem{patterson1979}
Patterson, J., Nather, R. E., Robinson, E. L., \& Handler, F. 1979.
\textit{Increase in the orbital period of AM Canum Venaticorum}.
ApJ, 232, 819.
\url{https://articles.adsabs.harvard.edu/pdf/1979ApJ...232..819P}.

\bibitem{roelofs2006}
Roelofs, G. H. A., Groot, P. J., Nelemans, G., Marsh, T. R., \& Steeghs, D. 2006.
\textit{Kinematics of the ultracompact helium accretor AM Canum Venaticorum}.
MNRAS, 371, 1231.
\url{https://academic.oup.com/mnras/article/371/3/1231/1007944}.

\bibitem{smak2023}
Smak, J. 2023.
\textit{AM CVn -- System Parameters and Gravitational Waves}.
Acta Astronomica, 73, 1.
\url{https://arxiv.org/pdf/2312.16506}.

\bibitem{solheim2010}
Solheim, J.-E. 2010.
\textit{AM CVn Stars: Status and Challenges}.
PASP, 122, 1133.

\end{thebibliography}

\appendix
\section{Individual Contributions}

TODO: Fill in team-member names and individual contributions before
submission, as required by the project instructions.

\section{Reproducibility}

The model suite is generated with:
\begin{verbatim}
python scripts/amcvn_physical_model_suite.py --bootstrap 250
\end{verbatim}
The generated CSV files use the \texttt{physical\_*} prefix in the AM CVn
periodogram folder.  The generated figures use the
\texttt{AM\_CVn\_physical\_*} prefix in the plot folder.

\end{document}
